% !TEX root = ../../../main.tex
% Week 4, IMG_9234, together with the bounded-sequence treatment in §1.9.
% Extended tail bounds are supporting notation for the ratio and root tests.
\section{Upper and lower limits (lim sup and lim inf)}
\label{sec:analysis-i:upper-and-lower-limits}

For a bounded real sequence $(a_n)$, define its tail bounds by
\[
  \alpha_n=\inf\{a_k\mid k\ge n\},\qquad
  \beta_n=\sup\{a_k\mid k\ge n\}.
\]
Removing early terms can only raise the infimum and lower the supremum.
Thus $(\alpha_n)$ is increasing, $(\beta_n)$ is decreasing, and
\[
  \liminf_{n\to\infty}a_n=\lim_{n\to\infty}\alpha_n,
  \qquad
  \limsup_{n\to\infty}a_n=\lim_{n\to\infty}\beta_n.
\]
By \autoref{thm:analysis-i:extreme-limit-points}, these are respectively
the smallest and largest limits of convergent subsequences. The full
sequence converges exactly when these two numbers agree. The tail
picture and its proof appear in \autoref{sec:analysis-i:denseness}.

\subsection{Allowing infinite tail bounds}

The tests for series also use sequences that need not be bounded.
We allow tail bounds to take values in $\R\cup\{-\infty,+\infty\}$.
An unbounded-above tail has supremum $+\infty$, and an
unbounded-below tail has infimum $-\infty$. The definitions become
\[
  \liminf a_n=\sup_n\alpha_n,\qquad
  \limsup a_n=\inf_n\beta_n.
\]
For bounded sequences these agree with the preceding definitions.
The infinite values record divergence of the tail bounds; they are
not additional real numbers.

\begin{marginfigure}
  \centering
  % !TEX root = ../../main.tex
% Shared physical dimensions for interval and number-line diagrams.
\tikzset{
  interval diagram/.style={
    x=1cm,y=6mm,font=\footnotesize,
    color=black,line width=0.5pt,line cap=round,line join=round,
    >={Stealth[length=4pt,width=3pt]}
  },
  interval endpoint/.style={
    circle,draw=black,line width=0.5pt,
    minimum size=4pt,inner sep=0pt,outer sep=0pt
  },
  interval open/.style={interval endpoint,fill=white},
  interval closed/.style={interval endpoint,fill=black},
  interval label/.style={inner sep=1pt}
}

\begin{tikzpicture}[interval diagram,x=3.8mm,y=2.6mm]
  \node at (4.5,9.8) {$a_n=n$};
  \draw[->] (0,0) -- (7.6,0) node[right=4pt] {$n$};
  \draw[->] (0,0) -- (0,8.1) node[above=5pt] {$a_n$};
  \fill[black!7] (3,2.5) rectangle (7.1,7.8);
  \draw[gray,dashed] (0,2.5) -- (7.1,2.5);
  \node[left=5pt] at (0,2.5) {$r$};
  \draw[gray,dashed] (3,0) -- (3,7.8);
  \node[below=6pt] at (3,0) {$N$};
  \foreach \n in {1,...,6}{
    \fill (\n,\n) circle[radius=1.4pt];
  }
  \node at (7.5,7.2) {$\cdots$};

  \begin{scope}[yshift=-33mm,x=3.8mm,y=1.5mm]
    \node at (4.5,11.4) {$b_n=(-1)^n n$};
    \draw[->] (0,0) -- (7.6,0) node[right=4pt] {$n$};
    \draw[->] (0,-7.5) -- (0,8.1) node[above=5pt] {$b_n$};
    \foreach \n in {1,...,6}{
      \pgfmathsetmacro{\term}{(-1)^\n*\n}
      \fill (\n,\term) circle[radius=1.4pt];
    }
    \node at (8,8.3) {$\cdots$};
    \node at (8,-8.3) {$\cdots$};
  \end{scope}
\end{tikzpicture}

  \caption[Two patterns of infinite tail bounds.]{For $a_n=n$, every sufficiently late term stays above a
    chosen level $r$. Here $r=5/2$ and $N=3$.
    For $b_n=(-1)^n n$, every tail reaches arbitrarily far in both
    directions. Dots are actual terms; ellipses indicate continuation.}
  \label{fig:infinite-tail-bounds}
\end{marginfigure}

\begin{example}[Two patterns of unbounded tails]
  \label{ex:analysis-i:infinite-tail-bounds}
  For $a_n=n$, the tail beginning at $N$ has $\alpha_N=N$ and
  $\beta_N=+\infty$. Hence both $\liminf a_n$ and $\limsup a_n$
  equal $+\infty$, and $a_n\to+\infty$.
  For $b_n=(-1)^n n$, every tail is unbounded above and below, so
  $\alpha_N=-\infty$ and $\beta_N=+\infty$ for every $N$.
  Its lower and upper limits are different, and it does not tend to
  either infinity. \autoref{fig:infinite-tail-bounds} compares the two
  patterns. Infinite tail bounds describe the whole tail, even though
  each individual term is a finite real number.
\end{example}

The feature we will need is an eventual bound. If $\limsup a_n<r$,
then $\beta_N<r$ for some $N$, so $a_n<r$ for every $n\ge N$.
Similarly, if $\liminf a_n>r$, then every sufficiently late term
exceeds $r$. The strict gap between the limit bound and $r$ is essential.

\begin{proposition}
  If $a_n\le b_n$ for every sufficiently large $n$, then
  $\liminf a_n\le\limsup b_n$.
\end{proposition}

\begin{proof}
  For all sufficiently late tails, $\alpha_n^{a}\le a_k\le b_k
  \le\beta_n^{b}$ whenever $k\ge n$. Given any two tail indices
  $i,j$, choose $k$ at least as large as both and beyond the index
  where $a_k\le b_k$. Then $\alpha_i^{a}\le\beta_j^{b}$.
  Taking the supremum over $i$ and then the infimum over $j$ gives
  the claimed inequality, including when a bound is infinite.
\end{proof}
